%% =====================================================================
%%  CO5151 · G4 · D3 FINAL REPORT (BẢN TIẾNG VIỆT — DRAFT CÓ SỐ LIỆU ĐÃ KIỂM)
%%  Đề tài: LegalPilot-VN — Hệ đa tác tử tư vấn tuân thủ pháp luật VN
%%  Hướng: RESEARCH (RT1–RT4)
%%
%%  BIÊN DỊCH (XeTeX):
%%    tectonic main.tex   → main.pdf
%%
%%  GIỚI HẠN: THÂN BÀI 6–8 TRANG (§1–§9). References + Appendix KHÔNG tính.
%%  Mọi \TODO{...} là phần chưa đo/chưa có từ nhóm — sẽ điền sau khi chạy xong.
%% =====================================================================
\documentclass[sigconf,nonacm]{acmart}

\usepackage[vietnamese]{babel}
\usepackage{booktabs}
\usepackage{graphicx}
\usepackage{xcolor}
\usepackage{multirow}
\usepackage{enumitem}
\usepackage{listings}
\usepackage{balance}
\usepackage{float}
\usepackage{tikz}
\usetikzlibrary{arrows.meta,positioning}

\newcommand{\TODO}[1]{\textcolor{red}{\textbf{[TODO: #1]}}}
\newcommand{\EVID}[1]{\textcolor{blue}{\footnotesize[#1]}}

\tolerance=1200
\setlength{\emergencystretch}{2.5em}

\setcopyright{none}
\settopmatter{printacmref=false}
\renewcommand\footnotetextcopyrightpermission[1]{}
\acmConference[CO5151 --- HK261]{Báo cáo cuối kỳ (Research Track: RT1--RT4)}{09/12/2026}{TP. Hồ Chí Minh, Việt Nam}

\begin{document}

\title{LegalPilot-VN: Hệ đa tác tử truy hồi tăng cường đồ thị tri thức cho tư vấn tuân thủ pháp luật Việt Nam}

\author{Đặng Lâm Tùng}
\affiliation{\institution{Trường ĐH Bách khoa --- ĐHQG-HCM}\country{Việt Nam}}
\email{tung.danglam@hcmut.edu.vn}

\author{Nguyễn Trung Phong}
\affiliation{\institution{Trường ĐH Bách khoa --- ĐHQG-HCM}\country{Việt Nam}}
\email{TODO-email-phong@hcmut.edu.vn}

\author{Vũ Việt Hưng}
\affiliation{\institution{Trường ĐH Bách khoa --- ĐHQG-HCM}\country{Việt Nam}}
\email{TODO-email-hung@hcmut.edu.vn}

\renewcommand{\shortauthors}{G4: Tùng, Phong, Hưng}

%% =====================================================================
\begin{abstract}
Chuyên viên tuân thủ ngân hàng Việt Nam phải tra cứu hàng nghìn văn bản quy phạm pháp luật (VBQPPL) có quan hệ sửa đổi, thay thế chằng chịt; chỉ một dẫn chiếu hết hiệu lực có thể kéo theo rủi ro pháp lý trực tiếp. Nhóm G4 trình bày \textbf{LegalPilot-VN} (hướng \textbf{Research, RT1--RT4}), một hệ đa tác tử kết hợp truy hồi lai (Qdrant + đồ thị Neo4j) với \emph{duyệt đồ thị chọn lọc theo hiệu lực}, \emph{kiểm toán từng mệnh đề} (\texttt{ClaimAuditorAgent}) và \emph{cổng kiểm chứng dẫn chiếu tất định} (\texttt{answer\_gate}) đặt ngoài prompt. Trên bộ 100 câu hỏi luật Ngân hàng Nhà nước (\texttt{sbv\_testset\_tvpl.json}, chấm 1 lượt trên cấu hình \texttt{p1}), hệ thống đạt \textbf{96/100} ở chỉ số proxy ba điều kiện định nghĩa bằng code (so với mức công bố \textbf{91/100} của SBV-LawGraph~\cite{phan2026sbv}), cùng độ chính xác dẫn chiếu 99,0\% và trúng căn cứ 97,0\%. Về bảo mật, chúng tôi xây dựng bộ 20 ca tấn công prompt injection gắn với điểm vào thật trong mã nguồn; đo sơ bộ 10 ca trước khi vá ghi nhận 6 ca bị chặn, 2 ca lọt ở mức unit/API (A16, A20), 1 ca chặn một phần (A07) và 1 ca không quy công được (A04). \emph{Giới hạn trọng yếu:} con số 96\% là \textbf{cận trên} vì chưa có LLM-judge chấm tương đương ngữ nghĩa như~\cite{phan2026sbv}, mới đo 1 lượt (chưa có trung bình $\pm$ SD qua 3 run), Precision@2 thực tế đạt 34,47\% (thấp hơn 37--39\% của paper và khác $k$), và 10 ca indirect còn ở trạng thái chờ chèn fixture. \textbf{Kho mã nguồn:} \url{https://github.com/CO5151/CO5151} $\cdot$ Tag: \texttt{final} $\cdot$ Commit: \TODO{điền commit hash 40 ký tự khi đóng gói D4}.
\end{abstract}

\maketitle

%% =====================================================================
\section{Giới thiệu}
\label{sec:intro}

\paragraph{Người dùng và bài toán.}
Đối tượng sử dụng mục tiêu của LegalPilot-VN là chuyên viên tuân thủ (compliance officer) và pháp chế doanh nghiệp trong lĩnh vực tài chính--ngân hàng tại Việt Nam. Khi thẩm định hồ sơ tín dụng, ví điện tử, hay khoản vay nước ngoài, người dùng phải đối chiếu trên kho 1.703 VBQPPL Ngân hàng Nhà nước (NHNN) với nhiều tầng quan hệ sửa đổi, hướng dẫn và thay thế (\texttt{SUPERSEDES}, \texttt{GUIDES}). Một sai sót thường gặp của các chatbot RAG thông thường là trộn lẫn văn bản còn hiệu lực với văn bản đã bị bãi bỏ (như Nghị định 101/2012/NĐ-CP sau khi Nghị định 52/2024/NĐ-CP có hiệu lực) hoặc bỏ sót điều khoản xếp hạng cao để tự suy diễn theo trí nhớ tham số.

\paragraph{Vì sao cần hệ đa tác tử thay cho RAG tĩnh.}
Truy vấn tuân thủ không thể giải quyết bằng một bước tìm kiếm $k$-NN cố định: câu hỏi thường chứa nhiều ý ràng buộc chéo đòi hỏi (i) phân rã thành các truy vấn con, (ii) mở rộng chọn lọc trên đồ thị pháp lý theo trạng thái hiệu lực, (iii) kiểm toán từng mệnh đề nguyên tử trong bản nháp đối chiếu với ngữ cảnh đã truy xuất, và (iv) quay lui loại bỏ văn bản hết hiệu lực hoặc bổ sung điều khoản còn thiếu khi phát hiện dẫn chiếu lệch. Mỗi bước phụ thuộc trực tiếp vào trạng thái trung gian của bước trước, tạo thành vòng phản hồi khép kín có điều kiện dừng rõ ràng.

\paragraph{Phạm vi và ngoài phạm vi.}
Phạm vi của hệ thống là hỗ trợ tra cứu và lập hồ sơ đối chiếu quy chuẩn (compliance dossier) trên kho luật NHNN với bộ đánh giá 100 câu hỏi có đáp án chuẩn chuyên gia. Hệ thống \textbf{không} cung cấp dịch vụ tư vấn lách luật, đại diện tranh tụng tại toà, hay tự động nộp hồ sơ hành chính khi chưa có phê duyệt xác thực của chuyên viên.

\paragraph{Đóng góp thực tế.}
Chỉ tính các thành phần đã cài đặt và có log kiểm chứng:
\begin{enumerate}[leftmargin=*]
  \item Cơ chế truy hồi lai (\texttt{LawGraphAgent}) kết hợp vector Qdrant và mở rộng đồ thị có lọc hiệu lực (\texttt{max\_hops=2}, \texttt{max\_traversal\_nodes=10}), thay cho duyệt 1-hop cố định của SBV-LawGraph~\cite{phan2026sbv}.
  \item Kiến trúc kiểm chứng hai tầng gồm \texttt{ClaimAuditorAgent} (kiểm mệnh đề + tra cứu \texttt{vbpl.vn}) và cổng dẫn chiếu tất định \texttt{answer\_gate.py} chạy ngay trong vòng lặp sửa lỗi của \texttt{LegalOrchestrator}.
  \item Quy tắc ràng buộc đầu ra bắt buộc bản nháp trích $\geq 1$ căn cứ trong top-3 đã truy xuất, giúp nâng tỷ lệ đạt đồng thời 3 điều kiện từ 91,0\% (\texttt{v1}/\texttt{v2}) lên 96,0\% (\texttt{p1}, 1 run).
  \item Bộ 20 ca kiểm thử prompt injection (7 direct, 12 indirect, 1 direct-API) gắn với điểm vào thật trong mã nguồn cùng script tái lập (\texttt{run\_cases.py}), phát hiện 5 lỗ hổng kỹ thuật (F1--F5) và 1 lỗi dữ liệu hiệu lực.
\end{enumerate}

\paragraph{Giả thuyết nghiên cứu có thể bác bỏ (RT1).}
\emph{Trên tập 100 câu \texttt{sbv\_testset\_tvpl.json} với cùng mô hình và ngân sách token, việc kết hợp duyệt đồ thị có chọn lọc và cổng dẫn chiếu tất định có ràng buộc top-3 (\texttt{C1}) đạt tỷ lệ thành công proxy 3 điều kiện cao hơn ít nhất $+4$ điểm phần trăm so với cấu hình loại bỏ duyệt đồ thị chọn lọc (\texttt{C2}) ở mức chênh lệch vượt quá 1 độ lệch chuẩn (SD) qua 3 lần chạy; giả thuyết bị \textbf{bác bỏ} nếu chênh lệch nằm trong 1 SD hoặc độ trễ trung bình tăng vượt $2\times$ mà không cải thiện tỷ lệ thành công.}

\begin{table}[t]
\caption{Phân định phần tự xây dựng và thành phần dùng lại.}
\label{tab:ownership}
\footnotesize
\begin{tabular}{@{}p{2.3cm}p{5.7cm}@{}}
\toprule
\textbf{Nhóm} & \textbf{Thành phần và vị trí trong kho mã nguồn} \\
\midrule
Nhóm tự xây dựng & \texttt{LegalOrchestrator}, \texttt{LawGraphAgent}, \texttt{ClaimAuditorAgent}, \texttt{DrafterAgent} (\texttt{src/agents/}); cổng \texttt{answer\_gate.py}, \texttt{guardrails.py}, \texttt{sanitizer.py}, \texttt{token\_gate.py} (\texttt{src/security/}); bộ công cụ pháp lý (\texttt{src/tools/}); bộ 20 ca tấn công (\texttt{final/security/attacks/}). \\
\midrule
Framework, model \& dữ liệu dùng lại & Google ADK (\texttt{ADKRunner}), LiteLLM, DeepSeek API (\texttt{deepseek-chat}), Gemini (\texttt{gemini-2.5-flash}), Qdrant, Neo4j, SQLite (WAL mode), bộ dữ liệu 1.703 VBQPPL \& 100 câu hỏi từ \texttt{sbv\_testset\_tvpl.json}~\cite{phan2026sbv}. \\
\bottomrule
\end{tabular}
\end{table}

%% =====================================================================
\section{Cơ sở lý thuyết và công trình liên quan}
\label{sec:related}

Điểm khác biệt của đề tài nằm ở \textbf{cơ chế kiểm chứng vòng lặp và cổng dẫn chiếu tất định bằng mã nguồn}, không dùng tên framework (Google ADK) hay số lượng tác tử làm bằng chứng về tính mới.

\paragraph{Baseline trực tiếp: SBV-LawGraph~\cite{phan2026sbv}.}
Phan et al.~\cite{phan2026sbv} xây dựng hệ Hybrid RAG + Knowledge Graph trên 1.703 văn bản NHNN và bộ 100 câu hỏi chuẩn, đạt Answer Correctness công bố 91\%. Tuy nhiên, tác giả tự ghi nhận hai khoảng trống: (i) duyệt đồ thị ở mức \emph{1-hop cố định}, dễ thu thập thiếu quan hệ chuỗi hoặc kéo thêm nhiễu; và (ii) chưa xử lý triệt để văn bản hết hiệu lực (vẫn gom cả NĐ 101/2012/NĐ-CP). Nhóm giải quyết phần (i) bằng \texttt{LawGraphAgent} và phần (ii) bằng bộ lọc trạng thái hiệu lực kết hợp đối chiếu trực tuyến \texttt{vbpl.vn} (dù mục~\ref{sec:security} cũng phát hiện dữ liệu nguồn vẫn còn bản ghi sai).

\paragraph{GraphRAG và truy hồi tăng cường phản chiếu.}
GraphRAG~\cite{edge2024graphrag} tổ chức tri thức theo cụm cộng đồng cho câu hỏi tổng hợp toàn cục, nhưng tốn chi phí dựng đồ thị và không mô hình hoá tính thời điểm hiệu lực của quy phạm pháp luật (tính \emph{non-monotonic}: luật mới bãi bỏ luật cũ). Các kiến trúc tác tử suy luận--hành động như ReAct~\cite{yao2023react}, Self-RAG~\cite{asai2024selfrag} và kiểm chứng mệnh đề dài SAFE~\cite{wei2024safe} đặt nền móng cho bước phân tách mệnh đề nguyên tử ở \texttt{ClaimAuditorAgent}, song các nghiên cứu này chưa tích hợp cổng trích dẫn tất định ở đầu ra hoặc đánh giá bảo mật theo điểm vào thực tế.

\paragraph{An toàn tác tử và Prompt Injection.}
Greshake et al.~\cite{greshake2023injection} chỉ ra rằng nội dung truy xuất từ dữ liệu ngoài là bề mặt tấn công \emph{indirect prompt injection} nguy hiểm cho ứng dụng LLM. AgentDojo~\cite{debenedetti2024injection} đánh giá động các tấn công trên tác tử dùng tool, còn Instruction Hierarchy~\cite{wallace2024instruction} đề xuất phân tầng chỉ thị ở mức huấn luyện. Trong thực tế triển khai trên khung có sẵn, nhóm tập trung vào kiểm soát ngoài prompt ở tầng kiến trúc và định lượng 20 ca tấn công sát với hệ thống LegalPilot-VN đang chạy.

%% =====================================================================
\section{Thiết kế hệ thống}
\label{sec:design}

\subsection{Kiến trúc và luồng điều khiển (R1)}
Luồng điều khiển thực tế cài đặt trong \texttt{LegalOrchestrator.run()} (\texttt{src/agents/orchestrator.py}) gồm bốn giai đoạn tuần tự có vòng quay lui khép kín (Hình~\ref{fig:arch}):
\begin{enumerate}[leftmargin=*]
  \item \textbf{Giai đoạn 0 --- Cổng đầu vào (\texttt{SecurityGuardrails}):} kiểm tra 7 mẫu regex đầu vào (Vector 2--8). Nếu dính Vector 4 (trích xuất prompt/key) hoặc Vector 5 (SQL injection), ném \texttt{SecurityViolationError} và dừng ngay. Nếu Vector 2 (lách luật), trả \texttt{refused}. Nếu Vector 8 (lệnh nộp hồ sơ), gắn cờ \texttt{human\_confirmation\_needed=True} (lưu ý: ở bản hiện tại hàm chưa \texttt{return} sớm — xem Mục~\ref{sec:security}).
  \item \textbf{Giai đoạn 1--2 --- Lập kế hoạch \& Truy hồi lai:} \texttt{plan\_subgoals()} tách câu hỏi thành 3--5 mục tiêu con và từ khoá pháp lý. \texttt{LawGraphAgent} truy vấn song song Qdrant (\texttt{top\_k=5}, \texttt{score\_threshold=0.65}) và đồ thị Neo4j (\texttt{max\_hops=2}, \texttt{max\_traversal\_nodes=10}). Nếu không thu được điều khoản nào, hệ từ bỏ an toàn (\texttt{insufficient\_evidence}, trả về \texttt{ABSTAIN\_MSG}).
  \item \textbf{Giai đoạn 3 --- Vòng lặp Drafter $\leftrightarrow$ Auditor $\leftrightarrow$ Citation Gate:} chạy tối đa \texttt{max\_retries = 3} lượt (tức tối đa 4 lần sinh nháp). Mỗi lượt, \texttt{DrafterAgent} tổng hợp hồ sơ tư vấn từ \texttt{active\_clauses}; \texttt{ClaimAuditorAgent} đối chiếu trạng thái hiệu lực; sau đó \texttt{gate\_answer()} (\texttt{answer\_gate.py}) kiểm tra tất định mọi dẫn chiếu \emph{``Điều N <số hiệu>''} so với tập \texttt{allowed\_from\_clauses(active\_clauses)}. Nếu thiếu dẫn chiếu thực tế có trong graph, hệ gọi \texttt{fetch\_articles()} để tự phục hồi ngữ cảnh; nếu vi phạm quy tắc top-3 (chỉ kiểm 1 lần) hoặc chứa văn bản hết hiệu lực, hệ loại bỏ văn bản vi phạm khỏi \texttt{active\_clauses} và quay lại sinh bản nháp mới.
  \item \textbf{Giai đoạn 4 --- Xuất bản có kiểm soát:} nếu hết vòng lặp mà cổng trích dẫn vẫn chặn, \texttt{strip\_citations()} bóc các trích dẫn chưa xác thực hoặc hệ từ chối trả lời khi không còn căn cứ hợp lệ (\texttt{ABSTAIN\_MSG}); trạng thái cuối được lưu vào \texttt{LegalAgentState}.
\end{enumerate}

\begin{figure}[t]
\centering
\begin{tikzpicture}[
  node distance=0.45cm and 0.45cm,
  box/.style={draw, rounded corners=2pt, font=\scriptsize, align=center, inner sep=3.5pt, fill=blue!4},
  gate/.style={draw, rounded corners=2pt, font=\scriptsize, align=center, inner sep=3.5pt, fill=orange!10},
  arr/.style={-{Stealth[scale=0.8]}, thick}
]
  \node[gate] (in) {User Query $\rightarrow$ \texttt{SecurityGuardrails}\\(V2--V8 regex gate)};
  \node[box, below=of in] (plan) {\texttt{Planner}: 3--5 subgoals $\rightarrow$ \texttt{LawGraphAgent}\\(Qdrant top-5 + Neo4j 2-hop, $\leq$10 node)};
  \node[box, below=of plan] (loop) {\texttt{DrafterAgent} $\leftrightarrow$ \texttt{ClaimAuditorAgent}\\(Backtrack loop, \texttt{max\_retries=3})};
  \node[gate, below=of loop] (out) {\texttt{answer\_gate.py} (Top-3 + Allowlist)\\Pass $\rightarrow$ Dossier \quad|\quad Fail $\rightarrow$ Retry / Abstain};
  \draw[arr] (in) -- (plan);
  \draw[arr] (plan) -- (loop);
  \draw[arr] (loop) -- (out);
  \draw[arr] (out.east) -- ++(0.45,0) |- node[pos=0.25, right, font=\tiny, align=left]{Retry\\$\leq 3$} (loop.east);
\end{tikzpicture}
\caption{Kiến trúc điều khiển thực tế trong \texttt{src/agents/orchestrator.py} (khớp 1--1 với mã nguồn).}
\label{fig:arch}
\end{figure}

\subsection{Tools và phân quyền (R2)}
Bảng~\ref{tab:tools} liệt kê các công cụ thực tế trong \texttt{src/tools/} cùng thuộc tính an toàn. Đáng chú ý, kiểm toán mã nguồn cho thấy \texttt{submit\_portal\_filing} là công cụ ghi duy nhất nhưng hiện \emph{chưa} gắn vào danh sách tool của ADK agent (\texttt{adk\_pipeline.py}), và hàm kiểm token mới dừng ở bước khác rỗng (\texttt{not approver\_token}).

\begin{table}[t]
\caption{Danh sách công cụ trong \texttt{src/tools/} và cơ chế bảo vệ.}
\label{tab:tools}
\scriptsize
\begin{tabular}{@{}p{2.35cm}ccp{2.65cm}@{}}
\toprule
\textbf{Công cụ} & \textbf{Đọc/Ghi} & \textbf{Token?} & \textbf{Timeout / Fallback} \\
\midrule
\texttt{query\_lawgraph} & Đọc & Không & Timeout cấu hình; fallback dense index nội bộ \\
\texttt{trace\_selective\_edge} & Đọc & Không & Giới hạn \texttt{max\_hops=2}, 10 node \\
\texttt{retrieve\_statutory\_\allowbreak provisions} & Đọc & Không & Trả rỗng $\rightarrow$ nhánh \texttt{ABSTAIN\_MSG} \\
\texttt{verify\_vbpl\_status} & Đọc & Không & \texttt{DOS\_QUERY\_TIMEOUT=12s}; fallback \texttt{statute\_cache} 24h \\
\texttt{search\_official\_\allowbreak legal\_web} & Đọc & Không & Allowlist domain cơ quan nhà nước \\
\texttt{submit\_portal\_filing} & \textbf{Ghi} & \textbf{Có*} & Chưa nối ADK; *chỉ kiểm token $\neq \emptyset$ (A20) \\
\bottomrule
\end{tabular}
\end{table}

\subsection{Quản lý context và bộ nhớ (R3)}
\begin{itemize}[leftmargin=*]
  \item \textbf{Bộ nhớ ngắn hạn:} trạng thái từng phiên lưu trong \texttt{LegalAgentState} (\texttt{src/memory/state\_models.py}), truyền có cấu trúc giữa các bước qua danh sách \texttt{active\_clauses} (đã cắt ngưỡng \texttt{top\_k=5}, chặn trần \texttt{max\_output\_tokens=2048}).
  \item \textbf{Bộ nhớ dài hạn và chính sách ghi/truy hồi:} \texttt{SQLiteMemoryManager} (\texttt{src/memory/sqlite\_manager.py}) quản lý 4 bảng trên \texttt{data/enterprise\_compliance.db} ở chế độ WAL với truy vấn tham số hoá: \texttt{enterprise\_profile}, \texttt{audit\_history}, \texttt{statute\_cache} (TTL 24 giờ cho tra cứu \texttt{vbpl.vn}) và \texttt{pending\_tokens}; ngoài ra có \texttt{data/adk\_sessions.db} lưu phiên ADK. \emph{Chính sách ghi/truy hồi hiện tại:} \texttt{statute\_cache} chỉ ghi sau khi parser \texttt{gazette\_tool} trả về trạng thái chuẩn hoá; tuy nhiên chưa có bộ lọc nội dung độc hại trước khi nạp hồ sơ doanh nghiệp hoặc lịch sử hỏi đáp xuyên phiên (bề mặt được kiểm thử ở ca A17).
\end{itemize}

\subsection{Ba trục nâng cao ngoài tools và memory (R4)}
\begin{enumerate}[leftmargin=*]
  \item \textbf{Truy hồi đồ thị chọn lọc theo hiệu lực (\texttt{LawGraphAgent}):} thay thế mở rộng 1-hop bằng mở rộng 2-hop có lọc cạnh \texttt{SUPERSEDES}/\texttt{GUIDES}. \emph{Phương án thay thế đã cân nhắc:} BFS toàn phần (gây bùng nổ ngữ cảnh, tăng độ nhiễu) và PageRank tĩnh (bỏ qua thứ tự hiệu lực theo thời gian).
  \item \textbf{Kiểm toán mệnh đề (\texttt{ClaimAuditorAgent}):} bóc tách kết luận pháp lý thành \texttt{AtomicClaim} và đối chiếu hiệu lực trực tiếp với đồ thị/VBPL để loại bỏ văn bản đã hết hiệu lực trong vòng quay lui. \emph{Thay thế:} dùng LLM-as-judge end-to-end (chi phí cao, không tất định, khó truy vết từng điều).
  \item \textbf{Cổng dẫn chiếu tất định (\texttt{answer\_gate.py}):} hàm thuần regex + chuẩn hoá mã văn bản kiểm 100\% trích dẫn xuất hiện trong nháp so với tập đã truy xuất, kết hợp ràng buộc top-3 trong vòng lặp. \emph{Thay thế đã thử:} siết allowlist quá chặt hoặc chặn cứng toàn câu ở D2 làm tỷ lệ từ chối tăng vọt (26,7\%, xem Mục~\ref{sec:limits}); đặt cổng vào trong vòng lặp sửa lỗi giữ được độ chính xác mà giảm tỷ lệ từ chối về 0\%.
\end{enumerate}

%% =====================================================================
\section{Triển khai}
\label{sec:impl}

\paragraph{Ngăn xếp công nghệ và giới hạn kỹ thuật.}
Hệ thống cài đặt bằng Python 3.13/3.14 với Google ADK, Pydantic Settings (\texttt{src/core/config.py}), cấu hình YAML (\texttt{configs/system\_config.yaml}), cơ sở dữ liệu Qdrant (chỉ mục vector), Neo4j (đồ thị tri thức) và SQLite WAL. Mô hình chính dùng cho thực nghiệm 100 câu là DeepSeek (\texttt{deepseek-chat}, \texttt{temperature=0.1}, \texttt{max\_output\_tokens=2048}), cùng cấu hình dự phòng \texttt{gemini-2.5-flash} và cục bộ \texttt{openai/llama3.2}. Lỗi ngoại lệ ở công cụ ngoài được bọc timeout (\texttt{timeout\_seconds=15}, \texttt{DOS\_QUERY\_TIMEOUT\_SECONDS=12}) và dự phòng bằng chỉ mục/cache cục bộ; tuy nhiên nhánh lỗi dispatch hiện chưa lọc văn bản thô từ thông điệp lỗi (ca A14).

\subsection{Trace thực thi đại diện (A2)}
Để minh hoạ luồng thật nối từ đầu vào đến đầu ra: trên ca kiểm thử trực tiếp \texttt{A06} (log lưu trong \texttt{final/security/attacks/results/before.jsonl}), (1) đầu vào đi qua \texttt{SecurityGuardrails.evaluate\_query()} và nhận phán quyết \texttt{allowed} (bộc lộ lỗ hổng F1 chưa gọi \texttt{InputSanitizer}); (2) \texttt{Plan} và \texttt{LawGraphAgent} truy xuất 6 điều khoản (NĐ 26/2025, NĐ 94/2025, TT 02/2025, TT 61/2024, TT 64/2024); (3) tại vòng lặp đầu (\texttt{attempts=1}), \texttt{answer\_gate} chặn nháp với lý do \emph{"có dẫn chiếu tới điều khoản không tồn tại"} và buộc quay lui sửa; (4) ở vòng kiểm lại, cổng xác nhận \emph{"mọi dẫn chiếu đều khớp điều khoản đã truy xuất"} (\texttt{guardrail\_verdicts[-1].ok=True}) và xuất ra dossier cuối cùng không làm theo chỉ thị ẩn. \TODO{Bổ sung một trace hoàn chỉnh của truy vấn nhiều chặng từ \texttt{eval/traces/} khi hoàn tất W1.5, kèm mã bước memory SQLite và thời gian từng công cụ.}

\subsection{Thay đổi so với báo cáo giữa kỳ D2}
Bảng~\ref{tab:d2_diff} tóm tắt ba điều chỉnh cốt lõi giữa D2 (25/10/2026) và hiện tại.

\begin{table}[t]
\caption{Thay đổi kỹ thuật và phương pháp so với D2.}
\label{tab:d2_diff}
\scriptsize
\begin{tabular}{@{}p{2.3cm}p{2.5cm}p{2.9cm}@{}}
\toprule
\textbf{Hạng mục} & \textbf{Ở mốc D2} & \textbf{Ở mốc D3 hiện tại} \\
\midrule
Cổng dẫn chiếu & Chỉ chạy cuối pipeline / chặn cứng & Chạy trong vòng lặp + tự nạp graph + ràng buộc top-3 (\texttt{p1}) \\
Kết luận Precision@2 & Ghi ``nhất quán'' với paper & Đính chính: 34,47\% < 37--39\%, khác $k$, chưa so sánh được \\
Đánh giá bảo mật & Liệt kê threat model v0 lý thuyết & Kiểm toán code tìm F1--F5 + bộ 20 ca có JSONL raw \\
\bottomrule
\end{tabular}
\end{table}

\subsection{Các phần chưa cài đặt}
Nhóm ghi nhận trung thực bốn phần chưa hoàn thiện ở bản hiện tại: (i) \textbf{LLM-as-judge} chấm tương đương ngữ nghĩa cho cấu hình \texttt{p1}; (ii) \textbf{OCR bảng biểu} cho tài liệu PDF ảnh scan; (iii) \textbf{công tắc thực thi thực sự} cho \texttt{enable\_selective\_traversal} (đang là cấu hình chết F3, cần sửa ở W1.1 trước khi chạy ablation \texttt{C2}); và (iv) \textbf{các bản vá bảo mật FIX-1..FIX-8} (đã lên kế hoạch ở \texttt{INJECTION\_PLAN.md}).

%% =====================================================================
\section{Đánh giá}
\label{sec:eval}

\subsection{Câu hỏi đánh giá (RQ)}
\begin{itemize}[leftmargin=*]
  \item \textbf{RQ1 (Hiệu quả tổng thể):} So với cấu hình nền và giá trị công bố của SBV-LawGraph~\cite{phan2026sbv}, hệ thống có cải thiện độ chính xác dẫn chiếu và khả năng trúng căn cứ trên 100 câu luật NHNN hay không?
  \item \textbf{RQ2 (Ablation):} Cơ chế ràng buộc trích dẫn top-3 và duyệt đồ thị chọn lọc đóng góp bao nhiêu vào mức tăng kết quả khi giữ nguyên mô hình và ngân sách?
  \item \textbf{RQ3 (Chi phí \& độ trễ):} Thời gian xử lý và chi phí token trên mỗi câu hỏi có nằm trong ngưỡng sử dụng thực tế của chuyên viên tuân thủ?
\end{itemize}

\subsection{Tập nhiệm vụ, tiêu chí và nguy cơ nhiễm dữ liệu}
Tập đánh giá gồm toàn bộ \textbf{100 câu hỏi} trong \texttt{data/benchmarks/sbv\_testset\_tvpl.json} (không lọc bỏ), trích từ văn bản công khai trên \texttt{thuvienphapluat.vn} kèm danh sách điều khoản chuẩn do chuyên gia gán nhãn. Thành công trên mỗi câu được chấm tất định bằng mã nguồn qua 3 điều kiện đồng thời: không từ chối trả lời (\texttt{Abstain = False}), có trích dẫn (\texttt{has\_citation}), mọi trích dẫn hợp lệ trong ngữ cảnh (\texttt{citations\_valid}), và trúng căn cứ chuẩn (\texttt{gt\_hit}).
\paragraph{Contamination.} Vì bộ 100 câu dựa trên dữ liệu pháp luật công khai và trùng với tập benchmark của~\cite{phan2026sbv}, tồn tại nguy cơ mô hình nền đã gặp văn bản khi tiền huấn luyện. Nhóm bổ sung bộ câu hỏi kiểm chứng độc lập (held-out) để đối chiếu \TODO{chèn kết quả 5 câu held-out tự soạn}.

\subsection{Cấu hình so sánh và kết quả chính}
Bảng~\ref{tab:main} tổng hợp dữ liệu đã đo được trực tiếp từ 3 tệp raw trong \texttt{data/benchmarks/} (\texttt{v1}, \texttt{v2}, \texttt{p1} — mỗi cấu hình 1 lượt chạy). So với \texttt{v1} và \texttt{v2} (cùng đạt 91,0\% ở proxy 3 điều kiện), cấu hình \texttt{p1} đạt \textbf{96/100 (96,0\%)}, tức cao hơn \textbf{+5,0 điểm tuyệt đối} so với con số 91,0\% của~\cite{phan2026sbv}.

\begin{table}[t]
\caption{Kết quả trên 100 câu \texttt{sbv\_testset\_tvpl.json} (chỉ số proxy 3 điều kiện không dùng judge; \texttt{v1, v2, p1} là 1-run baseline).}
\label{tab:main}
\scriptsize
\begin{tabular}{@{}lrrrrr@{}}
\toprule
\textbf{Chỉ số} & \textbf{v1} & \textbf{v2} & \textbf{C1 (\texttt{p1})} & \textbf{C2 (Abl.)} & \textbf{Paper~\cite{phan2026sbv}} \\
\midrule
Abstain             & 0,0\%  & 1,0\%  & 0,0\%  & \TODO{} & --- \\
Citation presence   & 100,0\% & 99,0\% & 100,0\% & \TODO{} & --- \\
Citation validity   & 93,0\% & 94,9\% & 99,0\% & \TODO{} & --- \\
Trúng căn cứ (\texttt{gt\_hit}) & 97,0\% & 95,0\% & 97,0\% & \TODO{} & --- \\
\textbf{Đạt cả 3 điều kiện} & \textbf{91,0\%} & \textbf{91,0\%} & \textbf{96,0\%} & \TODO{} & \textbf{91,0\%} \\
\midrule
Mean $\pm$ SD (3 runs) & --- & --- & \TODO{} & \TODO{} & --- \\
Tổng số retry (100 câu) & 29 & 37 & 30 & \TODO{} & --- \\
Thời gian batch (s) & 1856,7 & 2056,4 & 1182,6 & \TODO{} & --- \\
Cost / task (USD) & \TODO{} & \TODO{} & \TODO{} & \TODO{} & --- \\
\bottomrule
\end{tabular}
\end{table}

\paragraph{Cảnh báo về thang đo.} Con số \textbf{96,0\% là cận trên} và không chứng minh vượt Answer Correctness của~\cite{phan2026sbv}: số 91\% của paper có thêm bước LLM-judge (Gemini) chấm \emph{tương đương ngữ nghĩa}, một điều kiện chỉ có thể làm giảm hoặc giữ nguyên điểm. Ngoài ra, kết quả ở Bảng~\ref{tab:main} mới đo 1 lượt nên chưa loại trừ được nhiễu ngẫu nhiên của LLM.

\subsection{Chỉ số truy hồi và nghiên cứu cắt bỏ (Ablation)}
Bảng~\ref{tab:retrieval} đối chiếu tầng truy hồi so với~\cite{phan2026sbv}. Hệ thống vượt ở Recall, F2-score và giảm nhiễu ngữ cảnh; nhưng \textbf{Precision@2 đạt 34,47\%, thấp hơn mức 37--39\% của paper} do nhóm tách 3--5 sub-query và đánh giá trên $k \approx 3,25$ nên chưa cùng mẫu số để kết luận vượt. Đồng thời, chỉ số hiệu lực 100\% được đánh dấu cảnh báo vì phát hiện metadata \texttt{101/2012/ND-CP} bị gán sai ở đồ thị (Mục~\ref{sec:security}).

\begin{table}[t]
\caption{Đối chiếu tầng truy hồi so với SBV-LawGraph~\cite{phan2026sbv}.}
\label{tab:retrieval}
\scriptsize
\begin{tabular}{@{}lrrp{2.7cm}@{}}
\toprule
\textbf{Chỉ số} & \textbf{Paper~\cite{phan2026sbv}} & \textbf{G4} & \textbf{Kết luận trung thực} \\
\midrule
Recall@k & 76,0\% (@10) & \textbf{97,0\%} & Vượt $+21,0$ điểm \\
F2-score & 0,60--0,61 & \textbf{0,684} & Vượt \\
Độ nhiễu ngữ cảnh & $>60,0\%$ & \textbf{$<15,0\%$} & Giảm $>45$ điểm \\
Lọc hết hiệu lực & 0,0\% & 100,0\%* & *Phụ thuộc metadata (\S\ref{sec:security}) \\
Precision@2 & 37,0--39,0\% & \textbf{34,47\%} & \textbf{Thấp hơn} (khác $k \approx 3,25$) \\
\bottomrule
\end{tabular}
\end{table}

Về cơ chế nhân quả (RQ2): thống kê trên 100 câu cho thấy căn cứ đúng nằm trong top-2 đạt 77\%, top-3 đạt 87\%, nhưng khi chưa ép ràng buộc chỉ 80\% câu trả lời trích đúng nó. Thử nghiệm riêng việc ép trích $\geq 1$ căn cứ top-3 trên 30 câu đẩy \texttt{gt\_hit} từ 80\% $\to$ 93,3\% và giảm abstain từ 6,7\% $\to$ 3,3\%; quy tắc ``mỗi kết luận chỉ nêu một điều'' ở \texttt{p1} giảm lỗi trích thừa từ 5--12 điều/câu xuống mức giúp \texttt{citations\_valid} đạt 99,0\% trên 100 câu. Phần ablation \texttt{C1} vs \texttt{C2} (tắt Selective Traversal) và chi phí token từng câu sẽ được bổ sung sau khi chạy đủ 600 task-run \TODO{chèn bảng C1 vs C2 và token cost sau W1.5--W1.8}.

\subsection{Phân tích lỗi sơ bộ trên cấu hình \texttt{p1}}
Trong 100 câu ở cấu hình \texttt{p1}, chỉ có \textbf{4 ca thất bại} (tức 4,0\%, lấy trực tiếp từ \texttt{eval\_correctness\_100\_p1.json}):
\begin{itemize}[leftmargin=*]
  \item \textbf{Nhóm trượt căn cứ chuẩn (\texttt{gt\_hit=False}, 3/100 ca = 3\%):} mã câu \texttt{\#24} (status \texttt{unverified}, 1 retry), \texttt{\#55} (\texttt{unverified}, 0 retry) và \texttt{\#65} (\texttt{ok}, 1 retry) đều có dẫn chiếu hợp lệ trong ngữ cảnh nhưng trỏ sang điều khoản lân cận thay vì điều khoản ground truth.
  \item \textbf{Nhóm sai dẫn chiếu (\texttt{citations\_valid=False}, 1/100 ca = 1\%):} mã câu \texttt{\#66} (\texttt{unverified}, 1 retry) trúng ý căn cứ nhưng kèm trích dẫn không xác thực qua cổng.
\end{itemize}
\TODO{Phân loại đầy đủ trên 300 lượt chạy của C1 khi có log trace hoàn chỉnh.}

%% =====================================================================
\section{Bảo mật và an toàn}
\label{sec:security}

\subsection{Threat model gắn hệ thống và phát hiện từ mã nguồn}
Dựa trên kiến trúc ở Mục~\ref{sec:design}, hệ thống đối diện bốn điểm vào (entry points) thực tế: (i) \textbf{câu hỏi người dùng}, (ii) \textbf{chunk truy xuất từ Qdrant/Neo4j và tài liệu nạp}, (iii) \textbf{output/lỗi của tool ngoài} (\texttt{vbpl.vn}, \texttt{web\_search}), và (iv) \textbf{bộ nhớ dài hạn SQLite}. Kiểm toán mã nguồn phát hiện 5 lỗ hổng thật: \textbf{F1}: \texttt{InputSanitizer} (\texttt{sanitizer.py}) không được gọi ở bất kỳ nơi nào trong \texttt{src/}; \textbf{F2}: cổng bảo vệ chỉ kiểm \texttt{query} đầu vào, mọi chunk và tool output đi thẳng vào prompt; \textbf{F3}: cờ \texttt{enable\_selective\_traversal} chưa được đọc trong luồng chạy; \textbf{F4}: Vector 8 (\texttt{requires\_human\_token}) bật cờ nhưng không dừng sinh dossier; \textbf{F5}: regex \texttt{\_CITE} trong \texttt{answer\_gate.py} dùng \texttt{[\^{}\textbackslash n]} nên dẫn chiếu bị ngắt dòng làm cổng mù. Ngoài ra, trên \texttt{data/processed/neo4j\_graph\_export.json}, bản ghi \texttt{101/2012/ND-CP} vẫn ghi \texttt{status: "in\_force"}, khiến đối chứng \texttt{A04} cũng viện dẫn nhầm nghị định này và làm chỉ số hiệu lực 100\% ở D2 mang tính vòng tròn.

\subsection{Kết quả 20 ca tấn công Prompt Injection}
Nhóm xây dựng 20 ca gắn đúng file/dòng mã nguồn (\texttt{final/security/attacks/cases.json}) và thống nhất 4 nhãn đánh giá rõ ràng: \texttt{BLOCKED} (hệ chặn đúng), \texttt{ATTACK\_SUCCESS} (hành vi cấm xảy ra — \emph{lưu ý: ``thành công'' là của kẻ tấn công}), \texttt{PARTIAL} (gắn cờ nhưng vẫn sinh nội dung), \texttt{INCONCLUSIVE} (ca đối chứng cũng kích hoạt hành vi). Bảng~\ref{tab:sec} trình bày kết quả Pha A đo ngày 26/11 (1 lần chạy/ca, 10 ca direct/unit/API đã chạy và 10 ca indirect đang ở trạng thái \texttt{STAGED} chờ chèn fixture).

\begin{table}[t]
\caption{Kết quả kiểm thử 20 ca bảo mật Pha A (trước khi vá, commit \texttt{77fb226}/\texttt{b40ace3}, 1 lần chạy/ca). Cột Pha B điền sau FIX-1..FIX-8.}
\label{tab:sec}
\scriptsize
\begin{tabular}{@{}p{0.6cm}p{1.2cm}p{1.65cm}p{1.7cm}p{1.1cm}p{0.8cm}@{}}
\toprule
\textbf{Mã} & \textbf{Loại} & \textbf{Điểm vào / Phòng vệ} & \textbf{Pha A (Trước vá)} & \textbf{Pha B} & \textbf{Tấn công?} \\
\midrule
A01 & direct & query / V4 regex & \texttt{BLOCKED} (refused) & \TODO{} & KHÔNG \\
A02 & direct & query / không từ khoá & \texttt{BLOCKED} (LLM tự từ chối) & \TODO{} & KHÔNG \\
A03 & direct & query / V2 regex & \texttt{BLOCKED} (refused) & \TODO{} & KHÔNG \\
A04 & direct & query / V3 + metadata & \texttt{INCONCLUSIVE} (control cũng dính) & \TODO{} & KQCG \\
A05 & direct & query / V5 SQLi regex & \texttt{BLOCKED} (refused) & \TODO{} & KHÔNG \\
A06 & direct & query / \texttt{Sanitizer} chưa gọi & \texttt{BLOCKED} (lỗ hổng hở, k đổi) & \TODO{} & KHÔNG \\
A07 & direct & query / V8 chưa return & \texttt{PARTIAL} (cờ True + vẫn soạn) & \TODO{} & MỘT PHẦN \\
\midrule
A16 & indirect & output / \texttt{\_CITE [\^{}\textbackslash n]} & \texttt{ATTACK\_SUCCESS} (0 trích dẫn) & \TODO{} & \textbf{CÓ} \\
A18 & indirect & path / \texttt{sanitize\_export\_path} & \texttt{BLOCKED} (strip âm thầm) & \TODO{} & KHÔNG \\
A20 & dir-API & \texttt{submit\_portal\_\allowbreak filing} & \texttt{ATTACK\_SUCCESS} (token \texttt{"OK"}) & \TODO{} & \textbf{CÓ} \\
\midrule
A08--15, A17, A19 & indirect (10 ca) & Qdrant / tool / PDF / memory / tên file & \texttt{STAGED} (đã sinh 10 fixture, chưa chèn chạy live) & \TODO{} & CHỜ CHÈN \\
\bottomrule
\end{tabular}
\end{table}

\paragraph{Ba bài học phương pháp rút ra từ Pha A.}
(1) \emph{Dương tính giả do in lại câu hỏi:} vì dossier in nguyên văn \texttt{YÊU CẦU TRA CỨU}, lần chấm đầu chấm nhầm A04/A06/A07; sau khi bổ sung hàm \texttt{strip\_echo()} trong \texttt{run\_cases.py}, A06 chuyển đúng về \texttt{BLOCKED}. (2) \emph{Thiếu kiểm soát $\neq$ tấn công thành công:} A06 chứng minh \texttt{InputSanitizer} bị bỏ qua (\texttt{status=allowed}), nhưng mô hình vẫn không thực thi payload override. (3) \emph{Một lần chạy chưa đủ:} ca đối chứng của A04 ở lần 1 không gọi NĐ 101/2012 nhưng lần 2 lại gọi, buộc Pha A đầy đủ và Pha B phải chạy $\geq 3$ lần/ca.

\subsection{Kiểm soát quyền, audit trail và rủi ro còn lại}
Hiện tại quyền ghi thực sự (\texttt{submit\_portal\_filing}) đã tách thành hàm riêng có kiểm tra cấu trúc và sinh mã giao dịch SHA-256, nhưng thiếu cơ chế xác thực chữ ký token thật (lỗ hổng A20) và thiếu log append-only nối toàn bộ chuỗi quyết định (FIX-8). \TODO{Chèn log audit trail thật của hành động A07 sau khi cài FIX-8.} Ba rủi ro còn lại sau khi vá dự kiến gồm: (a) tấn công diễn giải ngữ nghĩa hoàn toàn không chứa từ khoá cấm; (b) chỉ thị ẩn nhúng dưới dạng ảnh trong PDF scan mà bộ tách text không đọc được; và (c) văn bản pháp lý hợp lệ được sắp đặt lệch ngữ cảnh gây hiểu nhầm cho người đọc.

%% =====================================================================
\section{Giới hạn, đe doạ tính hợp lệ và khả năng mở rộng}
\label{sec:limits}

\begin{itemize}[leftmargin=*]
  \item \textbf{Giới hạn thiết kế:} chưa tích hợp LLM-as-judge kiểm tra tương đương ngữ nghĩa trên bản \texttt{p1}; chưa hỗ trợ OCR cho văn bản PDF dạng ảnh scan; công cụ ghi \texttt{submit\_portal\_filing} mới ở mức mô phỏng cổng dịch vụ công và chưa có xác thực token chữ ký số (FIX-8). Ngoài ra, thí nghiệm thất bại ở D2 khi áp dụng \emph{cổng chặn cứng hoàn toàn} không cho quay lui đã đẩy tỷ lệ từ chối lên 26,7\%, cho thấy cần cân bằng giữa an toàn trích dẫn và tính hữu dụng.
  \item \textbf{Đe doạ tính hợp lệ của phép đo (Threats to Validity):} (i) số 96/100 là 1 lượt đo và thiếu semantic equivalence; (ii) tập 100 câu là mẫu nhỏ trên nguồn công khai của một lĩnh vực duy nhất (NHNN), không thể suy rộng ra toàn bộ hệ thống pháp luật; (iii) chỉ số \texttt{Precision@2} chưa đo đồng nhất $k$ với paper; (iv) lỗ hổng dữ liệu \texttt{101/2012/ND-CP} làm phép đo hiệu lực mang tính vòng tròn với metadata đồ thị.
  \item \textbf{Khả năng mở rộng:} kiến trúc tách tầng thu thập--kiểm toán--cổng dẫn chiếu cho phép mở rộng sang lĩnh vực thuế, lao động bằng cách nạp đồ thị và thay bộ luật allowlist mà không phải huấn luyện lại mô hình.
\end{itemize}

%% =====================================================================
\section{Đạo đức và sử dụng có trách nhiệm}
\label{sec:ethics}

Kho dữ liệu gồm 1.703 VBQPPL và bộ câu hỏi kiểm thử được thu thập từ nguồn pháp luật công khai của Nhà nước Việt Nam; việc kiểm thử bảo mật thực hiện hoàn toàn trên môi trường cục bộ và dữ liệu giả lập (mock VBPL/DVC), không gửi payload tới tài khoản hay hạ tầng của bên thứ ba. Hệ thống không lưu dữ liệu khách hàng thực tế (hồ sơ doanh nghiệp kiểm thử lưu cục bộ trong \texttt{workspace/}). Về trách nhiệm vận hành, LegalPilot-VN chỉ đóng vai trò công cụ hỗ trợ tra cứu có dẫn chiếu kiểm chứng, \textbf{không} thay thế kết luận pháp lý chính thức; mọi hành động ghi không thể hoàn tác đều yêu cầu phê duyệt và token xác thực của chuyên viên có thẩm quyền.

%% =====================================================================
\section{Kết luận}
\label{sec:concl}

LegalPilot-VN chứng minh được hướng tiếp cận kết hợp truy hồi đồ thị chọn lọc, kiểm toán mệnh đề và cổng dẫn chiếu tất định có ràng buộc top-3 đạt mức thực nghiệm sơ bộ 96/100 trên proxy ba điều kiện (cao hơn mốc công bố 91/100 của SBV-LawGraph~\cite{phan2026sbv}, nhưng là cận trên vì thiếu semantic judge) cùng tỷ lệ hợp lệ dẫn chiếu 99,0\%. Đồng thời, quá trình kiểm toán mã nguồn và 20 ca tấn công đã chỉ ra chính xác các điểm mù ở dữ liệu truy xuất, regex ngắt dòng (A16), cơ chế human token (A20) và metadata hết hiệu lực (\texttt{101/2012}). \TODO{Chốt kết luận bác bỏ hay không bác bỏ giả thuyết RT1 kèm Mean $\pm$ SD qua 3 lần chạy của C1 vs C2 sau khi hoàn tất 600 task-run.}

%% =====================================================================
\balance
\bibliographystyle{ACM-Reference-Format}
\bibliography{refs}

%% =====================================================================
\appendix

\section{Tái lập và thông tin bàn giao D4}
\label{app:repro}

Bảng~\ref{tab:d4} tổng hợp đầy đủ 12 trường bàn giao theo §8 của \texttt{03-Bieu-mau-Final.pdf}. Ghi nhận một lỗi tái lập hiện có: môi trường \texttt{.venv/} mặc định thiếu gói \texttt{pydantic\_settings} và \texttt{bs4}; cần chạy bằng \texttt{.venv-adk/bin/python3} cho tới khi đóng gói \texttt{requirements.lock.txt} (W4.1).

\begin{table}[H]
\caption{Thông tin bàn giao D4 theo §8 biểu mẫu final.}
\label{tab:d4}
\scriptsize
\begin{tabular}{@{}p{2.3cm}p{5.6cm}@{}}
\toprule
\textbf{Hạng mục} & \textbf{Giá trị / Trạng thái} \\
\midrule
Repo URL & \url{https://github.com/CO5151/CO5151} \\
Tag / Commit hash & Tag: \texttt{final} $\cdot$ Hash: \TODO{40-char hash bản nộp} \\
Lệnh cài đặt & \TODO{Cập nhật sau W4.1: pip install -r requirements.lock.txt} \\
Lệnh chạy demo & \TODO{Cập nhật sau W4.4: make demo} \\
Lệnh đánh giá & \texttt{.venv-adk/bin/python3 scripts/eval\_correctness.py} · \texttt{run\_cases.py} \\
Lệnh tạo bảng & \TODO{Cập nhật sau W1.6: make tables} \\
Dữ liệu \& truy cập & \texttt{data/benchmarks/sbv\_testset\_tvpl.json}; kiểm quyền GVHD (W4.7) \\
Model/version, ngày & \texttt{deepseek-chat} (\texttt{p1} đo 11/2026); red-team Pha A 26/11/2026 \\
Hardware & \TODO{CPU/RAM máy chạy 600 task-run} \\
Thời gian / Chi phí & 100 câu \texttt{p1}: 1182,6s ($\approx 11,8$s/câu); chi phí USD: \TODO{} \\
Video dự phòng & \TODO{Link video ~10' quay trên tag final (W5.3)} \\
Raw results / logs & \texttt{data/benchmarks/eval\_correctness\_100*.json}; \texttt{final/security/attacks/results/} \\
Giới hạn tái lập & Cần API key + Neo4j/Qdrant cho chạy live; \texttt{.venv} thiếu 2 gói (W4.1) \\
\bottomrule
\end{tabular}
\end{table}

\section{Bảng phản hồi góp ý Proposal / D2}
\label{app:response}

Bảng~\ref{tab:feedback} trình bày các góp ý đã xử lý từ vòng D1 (theo §1 của \texttt{03-Bieu-mau-Final.pdf}) và chỗ chờ bổ sung góp ý bằng văn bản của D2.

\begin{table}[H]
\caption{Phản hồi góp ý Proposal / D2 (mẫu §1).}
\label{tab:feedback}
\scriptsize
\begin{tabular}{@{}p{0.5cm}p{1.9cm}p{2.0cm}p{2.1cm}p{0.9cm}@{}}
\toprule
\textbf{ID} & \textbf{Góp ý \& nguồn} & \textbf{Hành động} & \textbf{Vị trí + bằng chứng} & \textbf{Trạng thái} \\
\midrule
R1 & Xử lý non-monotonicity \& bẫy luật cũ (D1) & Cài quan hệ \texttt{SUPERSEDES} + Auditor + cổng trích dẫn & §3.4, §5.4 (\texttt{citations\_valid=99\%}) & Xong (lưu ý metadata 101/2012) \\
R2 & Kiểm soát chi phí token \& độ trễ (D1) & Đặt \texttt{max\_retries=3}, \texttt{top\_k=5}, cache VBPL 24h & §3.1, §3.3 (\texttt{p1} giảm 1856s $\to$ 1182s) & Một phần (thiếu cost USD) \\
R3 & Tính tái lập \& khách quan đánh giá (D1) & Định nghĩa proxy tất định bằng code + tách harness 20 ca & §5.2, \texttt{run\_cases.py}, \texttt{answer\_gate.py} & Một phần (chờ 3 runs) \\
R4 & \TODO{Góp ý D2 từ GVHD} & \TODO{Cập nhật sau W6.1} & \TODO{} & Chưa \\
\bottomrule
\end{tabular}
\end{table}

\section{Đóng góp thành viên và tuyên bố sử dụng AI}
\label{app:contrib}

Bảng~\ref{tab:contrib} tuân thủ đủ 5 cột theo mẫu §5 của \texttt{03-Bieu-mau-Final.pdf}. Các thành viên sẽ gửi đánh giá nội bộ cá nhân riêng qua email cho giảng viên trước hạn 06/12/2026.

\begin{table}[H]
\caption{Bảng phân công và đóng góp thực tế của nhóm G4.}
\label{tab:contrib}
\scriptsize
\begin{tabular}{@{}p{1.8cm}p{1.6cm}p{1.8cm}p{1.6cm}p{0.6cm}@{}}
\toprule
\textbf{Thành viên / MSHV} & \textbf{Phụ trách} & \textbf{Sản phẩm thực tế} & \textbf{Bằng chứng} & \textbf{Tỷ lệ} \\
\midrule
Đặng Lâm Tùng \TODO{MSHV} & Kiến trúc, ADK, eval harness, báo cáo & Orchestrator, cổng trích dẫn, khung D3/slide, 100-q eval & \texttt{b40ace3}, \texttt{77fb226}, §3, §5 & \TODO{}\% \\
Nguyễn Trung Phong \TODO{MSHV} & Ingestion, KG, memory, đóng gói D4 & Pipeline Neo4j/Qdrant, SQLite WAL, tool VBPL, D4 & \texttt{sqlite\_manager.py}, Phụ lục A & \TODO{}\% \\
Vũ Việt Hưng \TODO{MSHV} & Claim Auditor, Threat Model, Red team & ClaimAuditorAgent, suite 20 ca injection, phân tích lỗi & \texttt{468a329}, \texttt{cases.json}, §6 & \TODO{}\% \\
\midrule
\multicolumn{4}{r}{\textbf{Tổng cộng (phải bằng 100\%):}} & \textbf{100\%} \\
\bottomrule
\end{tabular}
\end{table}

\paragraph{Tuyên bố sử dụng AI (AI-use statement).}
Trong quá trình phát triển và biên soạn báo cáo, nhóm sử dụng \TODO{liệt kê tên công cụ/model AI + phiên bản đã dùng, ví dụ trợ lý lập trình IDE, Claude/Gemini/DeepSeek} cho các công đoạn hỗ trợ: gợi ý mã nguồn Python/LaTeX, sinh dữ liệu mẫu fixture cho bộ 20 ca kiểm thử bảo mật và đối chiếu ngữ pháp tiếng Việt/tiếng Anh. Nhóm tự thiết kế kiến trúc, trực tiếp viết/chạy các kịch bản đánh giá và kiểm chứng lại 100\% số liệu từ tệp JSON thô (\texttt{eval\_correctness\_100*.json}, \texttt{before.jsonl}), tự kiểm tra biên dịch mã nguồn và đọc chéo từng trích dẫn trong \texttt{refs.bib}.

\section{Bảng đối chiếu yêu cầu và bằng chứng}
\label{app:matrix}

Bảng~\ref{tab:req} đối chiếu các yêu cầu trong đề tài (R1--R4, R6 từ tài liệu Final/D1; R5, R7--R9 sẽ đối chiếu lại đề cương gốc khi hoàn tất) và yêu cầu hướng Research (RT1--RT4).

\begin{table}[H]
\caption{Ma trận đối chiếu yêu cầu và bằng chứng (mẫu §4).}
\label{tab:req}
\scriptsize
\begin{tabular}{@{}p{1.0cm}p{1.6cm}p{2.7cm}p{2.1cm}@{}}
\toprule
\textbf{Mã} & \textbf{Vị trí bài} & \textbf{Code / Config / Log} & \textbf{Trạng thái / Giới hạn} \\
\midrule
R1 & §\ref{sec:design}.1 & \texttt{src/agents/orchestrator.py} & Đạt (có Hình~\ref{fig:arch}) \\
R2 & §\ref{sec:design}.2 & \texttt{src/tools/*}, \texttt{public\_service\_portal.py} & Một phần (lỗ hổng A20, chưa nối ADK) \\
R3 & §\ref{sec:design}.3 & \texttt{src/memory/sqlite\_manager.py} & Một phần (chờ lọc ghi A17) \\
R4 & §\ref{sec:design}.4 & \texttt{lawgraph.py}, \texttt{claim\_auditor.py}, \texttt{answer\_gate.py} & Đạt (3 trục nâng cao) \\
R6 & §\ref{sec:security} & \texttt{guardrails.py}, \texttt{final/security/attacks/} & Một phần (10/20 ca đo Pha A) \\
\midrule
RT1 & §\ref{sec:intro}, §\ref{sec:concl} & Giả thuyết bác bỏ được (proxy + SD) & Đang chờ 3 runs C1/C2 \\
RT2 & §\ref{sec:related}, §\ref{sec:eval}.3 & Baseline \cite{phan2026sbv} + cần baseline chạy live & Hạn chế (khác judge, cần W1) \\
RT3 & §\ref{sec:eval}.4 & Can thiệp top-3 (30q) + \texttt{C1} vs \texttt{C2} & Một phần (chờ FIX-7 cho C2) \\
RT4 & §\ref{sec:eval}.2, Phụ lục A & \texttt{eval\_correctness\_100*.json}, \texttt{run\_cases.py} & Một phần (cần Mean$\pm$SD + \texttt{lock}) \\
\bottomrule
\end{tabular}
\end{table}

\end{document}
